% !TEX root = ../main.tex

\section{Test and Validation Protocol}
\label{sec:validation}

\subsection{Cross-dataset generalization test}
\label{subsec:crossdataset}

Both models are trained and early-stopped exclusively on SALICON splits (\secref{subsec:training11}); SALICON-only validation would measure in-distribution fit rather than genuine generalization. For a fair comparison we additionally evaluate both, for the first time, on MIT1003~\cite{judd2009learning}: 1{,}003 images with fixations from 15 observers, covering a substantially different distribution (general scenes, rather than SALICON's context-driven web imagery) never seen during training or model selection. The same preprocessing ($240\times320$ resize, no augmentation) and evaluation code are applied to both models on the same 1{,}003 images, so the comparison in \secref{sec:testing} isolates the encoder rather than pipeline differences.

\subsection{Evaluation metrics}
\label{subsec:metrics}

Saliency evaluation is sensitive to metric choice, since different metrics emphasize different map properties and a model can score well on one and poorly on another~\cite{bylinskii2019different}. We therefore report the same three complementary, distribution-based metrics used as loss terms in \secref{subsec:loss11}: KLD (\eq{eq:kld}, $\downarrow$), CC (\eq{eq:cc}, $\uparrow$) and SIM (\eq{eq:sim}, $\uparrow$), computed between the fused output $S$ and the MIT1003 ground truth, plus the total loss $\mathcal{L}_{\mathrm{total}}$ (\eq{eq:total}) used for checkpoint selection. Location-based metrics (AUC, NSS) are complementary but not yet part of our pipeline (\secref{sec:conclusions}).

\subsection{Model selection}
\label{subsec:selection}

For both models we test the checkpoint with the lowest SALICON validation loss observed before training stopped (\tab{tab:train11}). Neither run reached its early-stopping patience budget before being manually interrupted for practical (Colab session-length) reasons, so the numbers in \secref{sec:testing} reflect an as-yet incomplete training schedule and should be read as preliminary.
